Este capítulo descreve como a pesquisa será conduzida. A Seção~\ref{sec:desenho} apresenta o desenho da pesquisa. A Seção~\ref{sec:selecao} descreve a seleção e a extração do trecho de código estudado. A Seção~\ref{sec:pipeline} detalha o \textit{pipeline} de transformação apoiado por LLM. A Seção~\ref{sec:criterios} define os critérios de avaliação.

\section{Desenho da Pesquisa}
\label{sec:desenho}

Este trabalho é uma pesquisa aplicada, de abordagem qualitativa e quantitativa, conduzida como estudo de caso exploratório \cite{gil2002,prodanov2013}. A pesquisa é aplicada porque busca gerar conhecimento para a solução de um problema concreto: a modernização de uma aplicação \textit{low-code} em uso. É exploratória porque o tema, a extração de regras de negócio de fórmulas Power Fx com apoio de LLMs, ainda é pouco estudado, e o objetivo é compreender o fenômeno e levantar evidências iniciais, e não testar hipóteses de forma generalizável.

O estudo de caso foi escolhido por permitir investigar o problema em profundidade, no seu contexto real, quando os limites entre o fenômeno e o contexto não são claramente definidos \cite{gil2002}. O caso estudado é o sistema \textit{Eventos}, uma aplicação desenvolvida em Microsoft Power Apps e utilizada pela empresa em que a pesquisa é realizada. A unidade de análise é um trecho de lógica de negócio escrito em Power Fx, extraído desse sistema conforme os critérios da Seção~\ref{sec:selecao}.

A pesquisa segue quatro etapas: seleção e extração do trecho a ser estudado; aplicação do \textit{pipeline} de transformação via LLM; execução dos testes de equivalência; e validação qualitativa com \textit{stakeholders} da empresa. Os resultados de cada etapa serão apresentados no Capítulo~\ref{cap:resultados}.

\section{Seleção e Extração da Cobaia}
\label{sec:selecao}

Chama-se aqui de \textit{cobaia} o trecho de código Power Fx que será submetido ao processo de modernização. Como o sistema \textit{Eventos} possui muitas telas e fórmulas, não é viável, no escopo deste trabalho, migrá-lo integralmente. Por isso, será selecionado um trecho representativo do problema descrito na Seção~\ref{sec:lowcode}, de acordo com os seguintes critérios:

\begin{itemize}
    \item \textbf{relevância para o negócio:} o trecho deve implementar uma regra que os usuários do sistema reconheçam como importante, como validações, cálculos ou decisões de fluxo;
    \item \textbf{complexidade lógica:} o trecho deve conter condições aninhadas ou combinadas, operações sobre fontes de dados e mais de um caminho de execução;
    \item \textbf{acoplamento à interface:} a regra deve depender do estado de componentes visuais, de modo que sua extração exija separar a lógica da apresentação;
    \item \textbf{observabilidade:} deve ser possível executar o trecho no ambiente do Power Apps com entradas controladas, para que seu comportamento sirva de referência nos testes.
\end{itemize}

A escolha será feita em conjunto com os responsáveis pelo sistema. O trecho selecionado será exportado da aplicação, junto com as informações necessárias para compreendê-lo: os nomes e tipos dos componentes referenciados, a estrutura das fontes de dados utilizadas e as variáveis globais ou de contexto envolvidas.

Antes de qualquer envio a um serviço de LLM, o material passará por anonimização. Nomes de tabelas, colunas, componentes e telas que revelem informações da empresa serão substituídos por identificadores genéricos; dados pessoais, endereços de e-mail, URLs e identificadores internos serão removidos ou substituídos por valores fictícios. A substituição será feita de forma consistente, preservando a estrutura e os tipos das expressões, para que a lógica permaneça intacta. A tabela de correspondência entre os nomes originais e os anonimizados será mantida apenas localmente e não será enviada ao modelo nem publicada neste trabalho. Esse cuidado atende às políticas de segurança da informação da empresa e à Lei Geral de Proteção de Dados Pessoais (Lei nº 13.709/2018).

\section{Pipeline de Transformação via IA}
\label{sec:pipeline}

A transformação do trecho selecionado será feita por um \textit{pipeline} de engenharia de \textit{prompt} em três etapas sequenciais. A divisão em etapas segue a ideia de passos pequenos e verificáveis defendida por \citeonline{fowler2018refactoring}: em vez de pedir ao modelo uma conversão direta, cada etapa produz um artefato intermediário que pode ser revisado antes de servir de entrada para a etapa seguinte. A Figura~\ref{fig:pipeline} resume o fluxo.

\begin{figure}[H]
    \centering
    \caption{\textit{Pipeline} de transformação via LLM.}
    \label{fig:pipeline}
    \small
    \begin{tabular}{c}
        \fbox{\parbox{0.6\textwidth}{\centering Trecho Power Fx anonimizado + contexto}} \\
        $\downarrow$ \\
        \fbox{\parbox{0.6\textwidth}{\centering \textbf{Etapa 1:} Documentação Semântica}} \\
        $\downarrow$ \textit{revisão humana} \\
        \fbox{\parbox{0.6\textwidth}{\centering \textbf{Etapa 2:} Tradução para Python}} \\
        $\downarrow$ \\
        \fbox{\parbox{0.6\textwidth}{\centering \textbf{Etapa 3:} Geração de Testes em Pytest}} \\
        $\downarrow$ \\
        \fbox{\parbox{0.6\textwidth}{\centering Avaliação (Seção~\ref{sec:criterios})}} \\
    \end{tabular}
    \legend{Fonte: elaboração da autora.}
\end{figure}

\subsection{Etapa 1: Documentação Semântica}

Na primeira etapa, o modelo recebe o trecho Power Fx anonimizado e o seu contexto (componentes, fontes de dados e variáveis) e é instruído a descrever, em linguagem natural e em formato estruturado, as regras de negócio implementadas. A descrição deve identificar as entradas de cada regra, as condições avaliadas, os resultados possíveis e os efeitos colaterais, como gravação de dados, notificações e navegação entre telas.

O objetivo é realizar a engenharia reversa semântica discutida por \citeonline{pereira2026modernizacao}: recuperar o que a fórmula faz, separado de como ela está escrita. O documento gerado será revisado pela autora, comparando cada regra descrita com o código original. Regras omitidas, incorretas ou inventadas pelo modelo serão registradas, pois indicam a qualidade da etapa, e corrigidas antes de seguir. O documento revisado passa a ser a especificação de referência para as etapas seguintes.

\subsection{Etapa 2: Tradução para Python}

Na segunda etapa, o modelo recebe o trecho original e a documentação revisada e é instruído a implementar as regras em Python, respeitando as seguintes restrições:

\begin{itemize}
    \item as regras de negócio devem ser implementadas como funções que não dependam de componentes visuais, recebendo como parâmetros os valores que, no Power Apps, vinham da interface;
    \item o acesso a dados deve ficar separado das regras, por meio de funções ou classes específicas, de modo que as regras possam ser executadas sem conexão com a fonte de dados;
    \item as reações de interface (\texttt{Notify}, \texttt{Navigate}) devem ser substituídas por valores de retorno que indiquem o resultado da decisão;
    \item o código deve usar anotações de tipo e nomes coerentes com a documentação da Etapa 1.
\end{itemize}

O resultado é um \textit{backend} desacoplado, no sentido discutido na Seção~\ref{sec:refatoracao}: a lógica fica isolada e pode ser executada e testada sem a interface.

\subsection{Etapa 3: Geração de Testes em Pytest}

Na terceira etapa, o modelo recebe a documentação revisada e o código Python e é instruído a gerar testes unitários com o \textit{framework} Pytest. Cada regra documentada deve ter ao menos um teste para cada caminho de execução, incluindo valores-limite (por exemplo, a capacidade exatamente atingida) e entradas em branco ou inválidas, que concentram as diferenças de semântica entre Power Fx e Python apontadas na Seção~\ref{sec:ia}.

Há, nessa etapa, um risco de circularidade: se o mesmo modelo gera o código e os resultados esperados dos testes, um erro de interpretação pode aparecer em ambos e os testes passariam mesmo com o comportamento incorreto. Para evitar isso, os testes gerados pelo modelo serão usados apenas para definir os \textit{casos} (as combinações de entrada). Os resultados esperados serão obtidos executando-se o trecho original no Power Apps, em ambiente de teste, com as mesmas entradas. Dessa forma, o sistema original funciona como oráculo, à semelhança dos testes de caracterização propostos por \citeonline{feathers2004working}.

\subsection{Registro do experimento}

Para permitir a reprodução e a análise do processo, serão registrados: o modelo de linguagem utilizado e sua versão; os parâmetros de geração; a data de execução; o texto integral de cada \textit{prompt}; as respostas obtidas; e as intervenções manuais realizadas entre as etapas. Os \textit{prompts} serão apresentados em apêndice.

\section{Critérios de Avaliação}
\label{sec:criterios}

A avaliação combina uma dimensão quantitativa, voltada à equivalência funcional, e uma dimensão qualitativa, voltada à utilidade do resultado para a empresa.

\subsection{Teste de equivalência funcional}

A implementação em Python será considerada equivalente ao trecho original, para os casos avaliados, quando produzir o mesmo resultado que o Power Apps para todas as entradas do conjunto de testes. Esse critério segue o adotado por \citeonline{pan2024lost}: código que apenas executa sem erro não é suficiente; ele precisa reproduzir o comportamento do original. Serão analisados os seguintes indicadores:

\begin{itemize}
    \item \textbf{taxa de aprovação:} proporção de casos de teste em que o resultado do Python coincide com o do Power Apps;
    \item \textbf{cobertura de regras:} proporção das regras presentes no código original que foram identificadas na documentação da Etapa 1;
    \item \textbf{cobertura de código:} proporção das linhas e ramificações do código Python exercitadas pelos testes;
    \item \textbf{intervenções manuais:} número e natureza das correções feitas pela autora em cada etapa do \textit{pipeline}.
\end{itemize}

Cada divergência encontrada será analisada para identificar sua origem (erro de interpretação na documentação, erro de tradução ou diferença de semântica entre as linguagens). Ressalta-se que os testes mostram a equivalência apenas para os casos executados e não provam a equivalência para todas as entradas possíveis; por isso, a escolha dos casos, com atenção a caminhos de execução e valores-limite, é parte do método.

\subsection{Validação qualitativa com \textit{stakeholders}}

A segunda dimensão avalia se a documentação e o código gerados são úteis para quem conhece e mantém o sistema. Serão realizadas entrevistas semiestruturadas com \textit{stakeholders} da empresa, como os responsáveis pelo sistema \textit{Eventos} e usuários que conhecem as regras de negócio envolvidas. Os participantes analisarão a documentação gerada e a estrutura do código Python e responderão a questões sobre:

\begin{itemize}
    \item a fidelidade da documentação às regras de negócio que conhecem;
    \item a clareza da documentação em comparação com a leitura direta das fórmulas Power Fx;
    \item a percepção sobre a manutenibilidade da versão em Python;
    \item a viabilidade de adotar o processo na modernização de outras partes do sistema.
\end{itemize}

As respostas serão analisadas de forma qualitativa, buscando pontos de concordância, divergências e sugestões de melhoria.

\section{Limitações}

O estudo se limita a um único sistema e a um trecho selecionado dele, e seus resultados não podem ser generalizados para todas as aplicações Power Apps nem para outras plataformas \textit{low-code}. A migração da interface de usuário e a integração com fontes de dados de produção estão fora do escopo. O comportamento de LLMs pode variar entre versões e execuções, o que é mitigado, mas não eliminado, pelo registro detalhado descrito na Seção~\ref{sec:pipeline}. Por fim, a validação qualitativa depende de um número reduzido de participantes, todos vinculados à mesma organização.
